\section{Introduction}
\label{sec:introduction}

A homogeneous polynomial $p\in\R[x_1,\ldots,x_m]$ of degree $d\geq1$ is
\emph{hyperbolic with respect to} $e\in\R^m$ if $p(e)\ne0$ and every root of
$t\mapsto p(te-x)$ is real for every $x\in\R^m$.
Write these roots, with multiplicity, as $\lambda_1(x),\ldots,\lambda_d(x)$.
The closed hyperbolicity cone is
\[
 \Lambda_+(p,e)=\{x\in\R^m:\lambda_j(x)\geq0\text{ for }1\leq j\leq d\}.
\]
We write $\Sym^N$ for the real symmetric $N\times N$ matrices and
$A\succeq0$ (respectively, $A\succ0$) for positive semidefiniteness
(respectively, positive definiteness).
A cone is \emph{spectrahedral} if it has the form
\begin{equation}
 K=\{x\in\R^m:L(x)\succeq0\},\qquad
 L(x)=\sum_{i=1}^m x_iA_i,\quad A_i\in\Sym^N,
 \label{eq:spectrahedral}
\end{equation}
for some finite $N\geq1$.
The geometric Generalized Lax conjecture asserts that every hyperbolicity
cone is spectrahedral. We give a counterexample.

\subsection{The explicit polynomial}

All indices in $\{1,2,3\}$ in the following construction are cyclic.
For $z,y\in\R^3$, let
\begin{equation}
 b(z,y)=\sum_{i=1}^3 z_i^2y_i^2
       -2\sum_{1\leq i<j\leq3}z_iy_iz_jy_j
       +\sum_{i=1}^3z_i^2y_{i+1}^2.
 \label{eq:form}
\end{equation}
This is the coefficient-one Choi--Lam biquadratic
form~\cite[Section~4]{choi-lam-1977}.
Define $Q(y)\in\Sym^3$ by
\[
 Q(y)_{ii}=y_i^2+y_{i+1}^2,\qquad
 Q(y)_{ij}=-y_iy_j\quad(i\ne j),
\]
so that $b(z,y)=z^\top Q(y)z$.
For $a\in\R$, $r\in\R^3$ and $T\in\Sym^3$, define the linear map
$\Phi_y:\Sym^4\to\Sym^4$ by
\begin{equation}
 \Phi_y\!\begin{pmatrix}a&r^\top\\r&T\end{pmatrix}
 =\begin{pmatrix}
    \tr(Q(y)T)&-r^\top Q(y)\\
    -Q(y)r&aQ(y)
   \end{pmatrix}.
 \label{eq:phi}
\end{equation}
The dependence on $y$ is homogeneous quadratic. If $\adj X$ denotes the
adjugate of $X$, set
\begin{equation}
 p(X,Z,y)=\det\big((\det X)Z-\Phi_y(\adj X)\big),
 \qquad (X,Z,y)\in\Sym^4\times\Sym^4\times\R^3.
 \label{eq:polynomial}
\end{equation}
This formula is polynomial even when $X$ is singular. The ambient space
has dimension $10+10+3=23$.

\begin{theorem}\label{thm:main}
The polynomial \eqref{eq:polynomial} is homogeneous of degree $20$ and
hyperbolic with respect to $e=(I_4,I_4,0)$, with $p(e)=1$.
Its closed hyperbolicity cone $K=\Lambda_+(p,e)$ is not spectrahedral:
for every finite $N\geq1$ and every real linear map
$L:\Sym^4\times\Sym^4\times\R^3\to\Sym^N$,
\[
 K\ne\{(X,Z,y):L(X,Z,y)\succeq0\}.
\]
\end{theorem}

The conclusion concerns the cone itself, so it allows every possible
defining pencil and every finite matrix size. The compact formula for $p$
is convenient for the proof, but one determinant factor can be removed.

\begin{corollary}\label{cor:reduced}
The quotient $q=p/\det X$ extends to a homogeneous polynomial of degree
$16$ in the same $23$ variables. It is hyperbolic with respect to $e$,
with $q(e)=1$, and $\Lambda_+(q,e)=K$. In particular, its closed
hyperbolicity cone is not spectrahedral.
\end{corollary}

We prove the degree reduction in Section~\ref{sec:reduced}, after the
obstruction argument. We make no claim that the dimension or degree is
minimal. The real symmetric convention also covers Hermitian pencils:
for a Hermitian matrix $H=A+iB$, positivity is equivalent to positivity
of the real symmetric matrix $\left(\begin{smallmatrix}A&-B\\B&A\end{smallmatrix}\right)$.

\subsection{Historical context}

G\aa rding's theory establishes the convexity of hyperbolicity
cones~\cite{garding-1959}. The original Lax conjecture asks for a
definite determinantal representation of homogeneous hyperbolic polynomials in three
variables~\cite{lax-1958}. Helton and Vinnikov proved the corresponding
definite determinantal representation theorem~\cite{helton-vinnikov-2007};
Lewis, Parrilo, and Ramana established its equivalence with Lax's
formulation~\cite{lewis-parrilo-ramana-2005}. In higher dimension, Br\"and\'en
constructed a real-zero polynomial none of whose positive powers admits
a definite determinantal representation~\cite[Theorem~3.3]{branden-2011}.
That obstruction does not exclude a different determinant, with additional
factors, from defining the same cone. The geometric conjecture concerns this remaining
possibility. Kummer exhibited the distinction concretely: the cone of
the specialized V\'amos polynomial has a $7\times7$ representation,
although no positive power of that polynomial has a definite
determinantal representation~\cite[Section~3]{kummer-2016-vamos}.
Raghavendra, Ryder, Srivastava, and Weitz proved exponential
lower bounds on the matrix size of spectrahedral representations for
suitable hyperbolicity cones~\cite{raghavendra-ryder-srivastava-2018}.
Recent positive results include Netzer's theorem for
hyperbolic cubics in five variables~\cite[Theorem~3.9]{netzer-2026}.

Kummer and Netzer prove spectrahedrality for cones of strictly
hyperbolic polynomials~\cite[Theorem~1]{kummer-netzer-2026-strict}.
Here \emph{strictly hyperbolic} means that $p(te-x)$ has distinct roots
whenever $x\notin\R e$. The polynomial~\eqref{eq:polynomial} fails this
hypothesis: at $w=(I_4,0,0)\notin\R e$ it gives
\[
 p(te-w)=(t-1)^{16}t^4.
\]
For the reduced polynomial of Corollary~\ref{cor:reduced}, the same line
gives $q(te-w)=(t-1)^{12}t^4$. Thus neither defining polynomial satisfies
the strict-case hypothesis.

A \emph{spectrahedral shadow} is a linear image of a set defined by an
affine real symmetric linear matrix inequality.
Netzer and Sanyal proved that a hyperbolicity cone is a spectrahedral
shadow when every nonzero boundary point is a smooth point of the
defining polynomial~\cite[Theorem~1.1]{netzer-sanyal-2015}.
Our result concerns a direct representation in the original variables;
it makes no assertion about representations using auxiliary variables.

The scalar ingredient is the coefficient-one Choi--Lam biquadratic
form~\cite[Section~4]{choi-lam-1977}, which reduces the coefficient of
the extra cyclic squared monomials from two in Choi's earlier
nonnegative form~\cite{choi-1975} to one. Its inability to dominate a nonzero
bilinear square follows from the extremality proved in
\cite[Theorem~4.4]{choi-lam-1977}; we give a short proof of precisely this
property. It is called \emph{weak extremality} by Blekherman, Rai\c{t}\u{a},
Shankar, and Sinn~\cite[Definition~2.7]{blekherman-raita-shankar-sinn-2022},
whose characterization uses Taylor expansions of dominated squares at
zeros~\cite[Theorem~3.2]{blekherman-raita-shankar-sinn-2022}.
The correspondence between nonnegative biquadratic forms and
positive maps on real symmetric matrices is classical; see
\cite[Section~3]{ha-2013}. Positive maps have also been used to obtain
certificates of conic polynomial stability~\cite{dey-gardoll-theobald-2021}.
Here the construction turns the scalar obstruction into an obstruction
to every finite pencil describing one hyperbolicity cone. A related
general strategy appears in Scheiderer's local sums-of-squares
obstructions to spectrahedral shadows~\cite[Section~4]{scheiderer-2018}.
Our proof extracts bilinear square minorants from a differentiated
operator-norm identity; it does not invoke a criterion for shadows.

Gonz\'alez Nevado states a positive solution of the geometric conjecture
in~\cite[Theorem~53, pp.~16--17]{gonzalez-nevado-2026}.
Appendix~\ref{app:opposition} records a counterexample to the
path-containment assertion used in that argument. The construction
and proof in the present paper are independent of that comparison.


\subsection{Proof strategy}

The scalar obstruction is particularly rigid: $b$ is nonnegative and
nonzero, but every bilinear form $\ell$ satisfying $\ell(z,y)^2\leq b(z,y)$
for all $z,y$ is zero. Section~\ref{sec:form} gives a short proof using
zeros and curves along which $b$ vanishes to fourth order.

The map $\Phi_y$ translates this form into a matrix inequality. In
Section~\ref{sec:hyperbolic} we prove hyperbolicity and identify the slice
of $K$ with $X\succ0$ as
\[
 Z\succeq\Phi_y(X^{-1}).
\]
The remaining argument assumes an arbitrary pencil defining $K$ and
extracts a bilinear square dominated by $b$.
First, Section~\ref{sec:pencil} rescales the pencil along
$(X,s^2Z,sy)$ and obtains an equivalent block matrix inequality.
Section~\ref{sec:norm} then sends $X^{-1}$ and $Z^{-1}$ toward rank-one projections;
the resulting comparison yields an exact operator-norm identity for
kernel projections of the pencil blocks. These projections are
parametrized by the unit sphere in $\R^4$.
Finally, Section~\ref{sec:tangent} differentiates one kernel projection
on a neighborhood of constant rank in that sphere. At a suitable base
point, the three-dimensional tangent space can be identified with all
the $z$-variables, so this
local identity controls $b$ for every $z$ and $y$.
Each entry of the resulting matrix
is a bilinear form whose square is bounded by $b$, giving a contradiction.

The mechanism connecting the biquadratic form to the geometry of the cone
is the rescaling and subsequent first variation of these kernel
projections. It applies to an arbitrary hypothetical pencil without
comparing its determinant with $p$.
All matrix norms below are Euclidean operator norms, denoted by
$\norm{\,\cdot\,}_{\op}$; vector norms are Euclidean.
